% Part: many-valued-logic
% Chapter: three-valued-logics
% Section: lukasiewicz

\documentclass[../../../include/open-logic-section]{subfiles}

\begin{document}

\olfileid[pt]{mvl}{inf}{luk}

\olsection{Lógica de \L ukasiewicz}

\begin{editorial}
  Este é um breve ``esboço'' de uma seção sobre a lógica de infinitos
  valores de \L ukasiewicz.
\end{editorial}

\begin{defn}\ollabel{def:lukasiewicz} A lógica de infinitos valores de
\L ukasiewicz~$\LogLuk[\infty]$ é definida pela matriz:
\begin{enumerate}
  \item A linguagem proposicional padrão $\Lang L_0$ com
  $\lnot$, $\land$, $\lor$, $\lif$.
  \item O conjunto de valores de verdade $V_\infty$.
  \item $1$ é o único valor designado, isto é, $V^+ = \{1\}$.
  \item As funções de verdade são dadas pelas seguintes funções:
  \begin{align*}
    \tf{\lnot}[\LogLuk](x) & = 1 - x\\
    \tf{\land}[\LogLuk](x,y) & = \min(x,y)\\
    \tf{\lor}[\LogLuk](x,y) & = \max(x,y)\\
    \tf{\lif}[\LogLuk](x,y) & = \min(1,1-(x-y)) = \begin{cases}
      1 & \text{se } x \le y\\
      1-(x-y) & \text{caso contrário.}
    \end{cases}
    \end{align*}
\end{enumerate}
A lógica $m$-valorada de \L ukasiewicz é definida da mesma maneira,
exceto que $V = V_m$.
\end{defn}

\begin{prop}
  A lógica $\LogLuk[3]$ definida por \olref[thr][luk]{def:lukasiewicz}
  é a mesma que $\LogLuk[3]$ definida por \olref{def:lukasiewicz}.
\end{prop}

\begin{proof}
  Isso pode ser visto comparando as tabelas de verdade dos conectivos
  dadas em \olref[thr][luk]{def:lukasiewicz} com as tabelas de verdade
  determinadas pelas equações em \olref{def:lukasiewicz}:
  \begin{center}
    \begin{tabular}{c|c} 
      $\tf{\lnot}$ & \\ 
      \hline  
      $1$ & $0$ \\ 
      $1/2$ & $1/2$ \\
      $0$ & $1$ 
    \end{tabular}
    \quad
    \begin{tabular}{c|ccc} 
      $\tf{\land}[\LogLuk[3]]$ & $1$ & $1/2$ & $0$ \\ 
      \hline 
      $1$ & $1$ & $1/2$ & $0$ \\ 
      $1/2$ & $1/2$ & $1/2$ & $0$\\ 
      $0$ & $0$ & $0$ & $0$ 
    \end{tabular}
    \\[2ex]
    \begin{tabular}{c|ccc} 
      $\tf{\lor}[\LogLuk[3]]$ & $1$ & $1/2$ & $0$ \\ 
      \hline 
      $1$ & $1$ & $1$ & $1$ \\ 
      $1/2$ & $1$ & $1/2$ & $1/2$ \\
      $0$ & $1$ & $1/2$ & $0$ 
    \end{tabular}
    \quad
    \begin{tabular}{c|ccc} 
      $\tf{\lif}[\LogLuk[3]]$ & $1$ & $1/2$ & $0$ \\ 
      \hline 
      $1$ & $1$ & $1/2$ & $0$ \\ 
      $1/2$ & $1$ & $1$ & $1/2$  \\ 
      $0$ & $1$ & $1$ & $1$ 
    \end{tabular}
  \end{center} 
\end{proof}

\begin{prop}\ollabel{prop:luk-infty-m}
  Se $\Gamma \Entails[\LogLuk[\infty]] !B$, então $\Gamma
  \Entails[\LogLuk[m]] !B$ para todo~$m \ge 2$.
\end{prop}

\begin{proof}
  Exercício.
\end{proof}

\begin{prob}
  Demonstre \olref[mvl][inf][luk]{prop:luk-infty-m}.
\end{prob}

Na verdade, a recíproca também vale.

A lógica de infinitos valores de \L ukasiewicz é a lógica difusa mais
popular. Na literatura de lógica difusa, o condicional é muitas vezes
definido como $\lnot !A \lor !B$. O resultado seria uma lógica forte
de Kleene com infinitos valores.

\begin{prob}
  Mostre que $(p \lif q) \lor (q \lif p)$ é uma
  tautologia de~$\LogLuk[\infty]$.
\end{prob}

\end{document}
